% OpenStax Calculus Volume 3 -- Section 2.6 Exercises: Quadric Surfaces
% Exercises reproduced from OpenStax, licensed under CC BY 4.0.
% Source: https://openstax.org/books/calculus-volume-3/pages/2-6-quadric-surfaces
\documentclass{exercises}
\title{OpenStax Calculus Volume 3}
\subtitle{Section 2.6 Exercises: Quadric Surfaces}
\sourceurl{https://openstax.org/books/calculus-volume-3/pages/2-6-quadric-surfaces}
\author{}
\date{}

\begin{document}
\maketitle


For the following exercises, sketch and describe the cylindrical surface of the given equation.

\begin{enumerate}[start=1]
\item \techrequired $x^{2}+z^{2}=1$
\item \techrequired $x^{2}+y^{2}=9$
\item \techrequired $z=\cos(\frac{\pi}{2}+x)$
\item \techrequired $z=e^{x}$
\item \techrequired $z=9-y^{2}$
\item \techrequired $z=\ln(x)$
\end{enumerate}

For the following exercises, the graph of a quadric surface is given.


(a) Specify the name of the quadric surface.


(b) Determine the axis of the quadric surface.

\begin{enumerate}[resume]
\item \textit{[Figure: This figure is a surface inside of a box. Its cross section parallel to the y z plane would be an upside down parabola. The outside edges of the 3-dimensional box are scaled to represent the 3-dimensional coordinate system.]}
\item \textit{[Figure: This figure is a surface inside of a box. It is an elliptical cone. The outside edges of the 3-dimensional box are scaled to represent the 3-dimensional coordinate system.]}
\item \textit{[Figure: This figure is a surface in the 3-dimensional coordinate system. There are two conical shapes facing away from each other. They have the x axis through the center.]}
\item \textit{[Figure: This figure is a surface in the 3-dimensional coordinate system. It is a parabolic surface with the x axis through the center.]}
\end{enumerate}

For the following exercises, match the given quadric surface with its corresponding equation in standard form.


(a) $\frac{x^{2}}{4}+\frac{y^{2}}{9}-\frac{z^{2}}{12}=1$


(b) $\frac{x^{2}}{4}-\frac{y^{2}}{9}-\frac{z^{2}}{12}=1$


(c) $\frac{x^{2}}{4}+\frac{y^{2}}{9}+\frac{z^{2}}{12}=1$


(d) $z=4x^{2}+3y^{2}$


(e) $z=4x^{2}-y^{2}$


(f) $4x^{2}+y^{2}-z^{2}=0$

\begin{enumerate}[resume]
\item Hyperboloid of two sheets
\item Ellipsoid
\item Elliptic paraboloid
\item Hyperbolic paraboloid
\item Hyperboloid of one sheet
\item Elliptic cone
\end{enumerate}

For the following exercises, rewrite the given equation of the quadric surface in standard form. Identify the surface.

\begin{enumerate}[resume]
\item $-x^{2}+36y^{2}+36z^{2}=9$
\item $-4x^{2}+25y^{2}+z^{2}=100$
\item $-3x^{2}+5y^{2}-z^{2}=10$
\item $3x^{2}-y^{2}-6z^{2}=18$
\item $5y=x^{2}-z^{2}$
\item $8x^{2}-5y^{2}-10z=0$
\item $x^{2}+5y^{2}+3z^{2}-15=0$
\item $63x^{2}+7y^{2}+9z^{2}-63=0$
\item $x^{2}+5y^{2}-8z^{2}=0$
\item $5x^{2}-4y^{2}+20z^{2}=0$
\item $6x=3y^{2}+2z^{2}$
\item $49y=x^{2}+7z^{2}$
\end{enumerate}

For the following exercises, find the trace of the given quadric surface in the specified plane of coordinates and sketch it.

\begin{enumerate}[resume]
\item \techrequired $x^{2}+z^{2}+4y=0,z=0$
\item \techrequired $x^{2}+z^{2}+4y=0,x=0$
\item \techrequired $-4x^{2}+25y^{2}+z^{2}=100,x=0$
\item \techrequired $-4x^{2}+25y^{2}+z^{2}=100,y=0$
\item \techrequired $x^{2}+\frac{y^{2}}{4}+\frac{z^{2}}{100}=1,x=0$
\item \techrequired $x^{2}-y-z^{2}=1,y=0$
\item Use the graph of the given quadric surface to answer the questions. \\ \textit{[Figure: This figure is a surface inside of a box. It is an oval solid on its side. The outside edges of the 3-dimensional box are scaled to represent the 3-dimensional coordinate system.]}
\begin{enumerate}[label=(\alph*)]
  \item Specify the name of the quadric surface.
  \item Which of the equations---$16x^{2}+9y^{2}+36z^{2}=3600$, $9x^{2}+36y^{2}+16z^{2}=3600$, or $36x^{2}+9y^{2}+16z^{2}=3600$---corresponds to the graph?
  \item Use b. to write the equation of the quadric surface in standard form.
\end{enumerate}
\item Use the graph of the given quadric surface to answer the questions. \\ \textit{[Figure: This figure is a surface inside of a box. It is a parabolic solid opening up vertically. The outside edges of the 3-dimensional box are scaled to represent the 3-dimensional coordinate system.]}
\begin{enumerate}[label=(\alph*)]
  \item Specify the name of the quadric surface.
  \item Which of the equations---$36z=9x^{2}+y^{2}$, $9x^{2}+4y^{2}=36z$, or $-36z=-81x^{2}+4y^{2}$---corresponds to the graph above?
  \item Use b. to write the equation of the quadric surface in standard form.
\end{enumerate}
\end{enumerate}

For the following exercises, the equation of a quadric surface is given.


(a) Use the method of completing the square to write the equation in standard form.


(b) Identify the surface.

\begin{enumerate}[resume]
\item $x^{2}+2z^{2}+6x-8z+1=0$
\item $4x^{2}-y^{2}+z^{2}-8x+2y+2z+3=0$
\item $x^{2}+4y^{2}-4z^{2}-6x-16y-16z+5=0$
\item $x^{2}+z^{2}-4y+4=0$
\item $x^{2}+\frac{y^{2}}{4}-\frac{z^{2}}{3}+6x+9=0$
\item $x^{2}-y^{2}+z^{2}-12z+2x+37=0$
\item Write the standard form of the equation of the ellipsoid centered at the origin that passes through points $A(2,0,0),B(0,0,1)$, and $C(\frac{1}{2},\sqrt{11},\frac{1}{2})$.
\item Write the standard form of the equation of the ellipsoid centered at point $P(1,1,0)$ that passes through points $A(6,1,0),B(4,2,0)$ and $C(1,2,1)$.
\item Determine the intersection points of elliptic cone $x^{2}-y^{2}-z^{2}=0$ with the line of symmetric equations $\frac{x-1}{2}=\frac{y+1}{3}=z$.
\item Determine the intersection points of hyperbolic paraboloid $z=3x^{2}-2y^{2}$ with the line of parametric equations $x=3t,y=2t,z=19t$, where $t\in \mathbb{R}$.
\item Find an equation of the quadric surface with points $P(x,y,z)$ that are equidistant from point $Q(0,-1,0)$ and plane of equation $y=1$. Identify the surface.
\item Find an equation of the quadric surface with points $P(x,y,z)$ that are equidistant from point $Q(0,2,0)$ and plane of equation $y=-2$. Identify the surface.
\item If the surface of a parabolic reflector is described by equation $400z=x^{2}+y^{2}$, find the focal point of the reflector.
\item Consider the parabolic reflector described by equation $z=20x^{2}+20y^{2}$. Find its focal point.
\item Show that quadric surface $x^{2}+y^{2}+z^{2}+2xy+2xz+2yz+x+y+z=0$ reduces to two parallel planes.
\item Show that quadric surface $x^{2}+y^{2}+z^{2}-2xy-2xz+2yz-1=0$ reduces to two parallel planes.
\item \techrequired The intersection between cylinder ${(x-1)}^{2}+y^{2}=1$ and sphere $x^{2}+y^{2}+z^{2}=4$ is called a Viviani curve. \\ \textit{[Figure: This figure is a surface inside of a box. It is a sphere with a right circular cylinder through the sphere vertically. The outside edges of the 3-dimensional box are scaled to represent the 3-dimensional coordinate system.]}
\begin{enumerate}[label=(\alph*)]
  \item Solve the system consisting of the equations of the surfaces to find the equations of the intersection curve. (Hint: Find $x$ and $y$ in terms of $z$.)
  \item Use a computer algebra system (CAS) to visualize the intersection curve on sphere $x^{2}+y^{2}+z^{2}=4$.
\end{enumerate}
\item Hyperboloid of one sheet $25x^{2}+25y^{2}-z^{2}=25$ and elliptic cone $-25x^{2}+75y^{2}+z^{2}=0$ are represented in the following figure along with their intersection curves. Identify the intersection curves and find their equations (Hint: Find $y$ from the system consisting of the equations of the surfaces.) \\ \textit{[Figure: This figure is a surface inside of a box. It is a hyperbolic paraboloid with a hyperbola of two sheets intersecting. The outside edges of the 3-dimensional box are scaled to represent the 3-dimensional coordinate system.]}
\item \techrequired Use a CAS to create the intersection between cylinder $9x^{2}+4y^{2}=18$ and ellipsoid $36x^{2}+16y^{2}+9z^{2}=144$, and find the equations of the intersection curves.
\item \techrequired A spheroid is an ellipsoid with two equal semiaxes. For instance, the equation of a spheroid with the z-axis as its axis of symmetry is given by $\frac{x^{2}}{a^{2}}+\frac{y^{2}}{a^{2}}+\frac{z^{2}}{c^{2}}=1$, where $a$ and $c$ are positive real numbers. The spheroid is called oblate if $c<a$, and prolate for $c>a$.
\begin{enumerate}[label=(\alph*)]
  \item The eye cornea is approximated as a prolate spheroid with an axis that is the eye, where $a=8.7\,\text{mm}$ and $c=9.6\,\text{mm}$. Write the equation of the spheroid that models the cornea and sketch the surface.
  \item Give two examples of objects with prolate spheroid shapes.
\end{enumerate}
\item \techrequired In cartography, Earth is approximated by an oblate spheroid rather than a sphere. The radii at the equator and poles are approximately $3963$ mi and $3950$ mi, respectively.
\begin{enumerate}[label=(\alph*)]
  \item Write the equation in standard form of the ellipsoid that represents the shape of Earth. Assume the center of Earth is at the origin and that the trace formed by plane $z=0$ corresponds to the equator.
  \item Sketch the graph.
  \item Find an equation of the intersection curve of the surface with plane $z=1000$ that is parallel to the xy-plane. The intersection curve is called a parallel.
  \item Find an equation of the intersection curve of the surface with plane $x+y=0$ that passes through the z-axis. The intersection curve is called a meridian.
\end{enumerate}
\item \techrequired A set of buzzing stunt magnets (or ``rattlesnake eggs'') includes two sparkling, polished, superstrong spheroid-shaped magnets well-known for children's entertainment. Each magnet is $1.625$ in. long and $0.5$ in. wide at the middle. While tossing them into the air, they create a buzzing sound as they attract each other.
\begin{enumerate}[label=(\alph*)]
  \item Write the equation of the prolate spheroid centered at the origin that describes the shape of one of the magnets.
  \item Write the equations of the prolate spheroids that model the shape of the buzzing stunt magnets. Use a CAS to create the graphs.
\end{enumerate}
\item \techrequired A heart-shaped surface is given by equation ${(x^{2}+\frac{9}{4}y^{2}+z^{2}-1)}^{3}-x^{2}z^{3}-\frac{9}{80}y^{2}z^{3}=0$.
\begin{enumerate}[label=(\alph*)]
  \item Use a CAS to graph the surface that models this shape.
  \item Determine and sketch the trace of the heart-shaped surface on the xz-plane.
\end{enumerate}
\item \techrequired The ring torus symmetric about the z-axis is a special type of surface in topology and its equation is given by ${(x^{2}+y^{2}+z^{2}+R^{2}-r^{2})}^{2}=4R^{2}(x^{2}+y^{2})$, where $R>r>0$. The numbers $R$ and $r$ are called the major and minor radii, respectively, of the surface. The following figure shows a ring torus for which $R=2$ and $r=1$. \\ \textit{[Figure: This figure is a surface inside of a box. It is a torus, a doughnut shape. The outside edges of the 3-dimensional box are scaled to represent the 3-dimensional coordinate system.]}
\begin{enumerate}[label=(\alph*)]
  \item Write the equation of the ring torus with $R=2$ and $r=1$, and use a CAS to graph the surface. Compare the graph with the figure given.
  \item Determine the equation and sketch the trace of the ring torus from a. on the xy-plane.
  \item Give two examples of objects with ring torus shapes.
\end{enumerate}
\end{enumerate}

\end{document}
